\section{Paid inference, funding, and settlement}
\label{paid-asynchronous-jobs-and-on-chain-coordination}

The first commercial service is hosted \textit{A}: one admitted model returns
estimated AI share, AI-involvement confidence, and author
attribution within its qualified scope. Slop Ninja serves a published,
qualified artifact in its own infrastructure; \(R_0\) is such an artifact once
its public package is published
(Section~\ref{model-interfaces-and-starter-architectures}). Customers buy access through an API rather than
operating the artifact themselves. A later \textit{B} service would offer
bounded revision through private models hosted in miners' Cloudflare accounts,
under the conditions in Section~\ref{later-b}.

The coordinator opens paid inference only after deciding that measured quality
supports the advertised service. It publishes the accepted versions, supported
workloads, limitations, and activation decision. Emission eligibility alone
does not authorize a paid product. Each product needs its own cost,
availability, privacy, and delivery qualification; book-length inputs remain
outside an offer until those properties have been measured at that size.

Customer inference principal is converted to subnet alpha where needed and
burned. It pays neither a direct miner fee nor an operator margin. The economic
hypothesis is that customer demand purchases alpha where conversion is required
and removes alpha from supply through burns. Compared with retaining those
tokens, this may support alpha's market value for holders, including miners
and other emission recipients. It guarantees neither price appreciation nor
spendable proceeds sufficient to cover costs. This market-value channel is
distinct from a possible reward allocation for verified paid use, which
distributes emissions available under the chain's rules. Miners and the operator
must fund their bills while any emissions they earn remain subject to chain
allocation, market value, and liquidity.

The full-principal choice puts the entire inference price toward alpha burning,
purchasing alpha when conversion is required. The operator's intended upside comes
through its holdings and any owner or validator emissions it is entitled to
earn. This commits the treasury to funding operations and refund reserves
separately, accepting that exposure in pursuit of a benefit shared with other
alpha holders. The policy's commercial viability remains to be measured.

\subsection{Pricing and cost calibration}\label{market-pricing-in-alpha}

The operator quotes a fixed price for a bounded workload using a versioned
reference cost and multiplier:
\[
  P_{\mathrm{USD}} = M_{\mathrm{price}} C_{\mathrm{model}},\qquad
  C_{\mathrm{total}} = C_{\mathrm{model}} + C_{\mathrm{other}}.
\]
Here \(C_{\mathrm{model}}\) is the expected cost of all model invocations per
delivered job: execution of the \textit{A} artifact and, for later revision or a
later CFG input profile, ambiguity resolution, proposal, ranking, and required
semantic checks.
\(C_{\mathrm{other}}\) includes non-model parsing, orchestration, storage,
delivery, payment and chain fees, and expected refund overhead. Numeric-model
runtime is counted once. Training and mandatory benchmark research have
separate budgets.

An exploratory value is \(M_{\mathrm{price}}=25\). Detection and revision may need different
values. The multiplier sizes customer principal destined for burning relative
to reference model cost. Calibration must weigh the intended economic effect
against competitive pricing and customer demand. It does not set an operator
margin, fix the share of emissions allocated to inference, or guarantee a
miner's margin. Cost profiles bind model versions, workload limits,
maximum attempts, tariff assumptions, and measurement dates. Overruns cannot
increase an accepted quote. Repricing affects future reservations.

The initial detection hypothesis is \$0.02 per 1,000 words, with
\$0.01--\$0.03 as the first range to test. The following planning scenarios
are not measurements of a qualified production service.

\begin{center}
\small
\begin{tabularx}{\linewidth}{Lrrr}
\toprule
Service per 1,000 words & Model cost & Delivery subtotal & Price at \(M_{\mathrm{price}}=25\) \\
\midrule
Raw-text detection & \$0.0002875--\$0.0004250 & \$0.0004--\$0.0009 & \$0.007188--\$0.01063 \\
Bounded revision & \$0.008289 & \$0.0106--\$0.0112 & \$0.2072 \\
\bottomrule
\end{tabularx}
\end{center}

A launch detection job tokenizes the canonical text and runs the raw-text
\textit{A} artifact (Section~\ref{input-contract}); it has no grammar parse and
no language-model call. Its model cost is the October 29, 2026 Container
busy-time measurement, \$0.23--\$0.34 per 1,000 documents at the 2,048-token
limit, used as the bound for a 1,000-word document within that
limit~\cite{containerserving}. The bounded revision scenario, dated October 22,
2026, uses published tariffs and assumed workloads for later \textit{B}, which
parses its source~\cite{cfpricing}: 8,000 input and 1,000 output tokens of
batched ambiguity resolution on Llama 3.1 8B fast FP8, at \$0.045/\$0.384 per
million input/output tokens; three proposal calls totaling 12,000 input and
4,500 output tokens at the Mistral 7B v0.1 tariff of \$0.110/\$0.190; an
independent preservation check using 8,000 input and 600 output tokens at the
Llama 3.3 70B fast FP8 tariff of \$0.293/\$2.253; and \$0.00002 for numeric
scoring and ranking, an allocation that predates the Container measurement.
Both model costs include a 25\% failure/retry allowance. Prices use the
unrounded model-cost totals; model costs and prices are displayed to four
significant figures.

The delivery subtotal includes model cost and adds \$0.0001--\$0.0005 per
detection or \$0.0002--\$0.0008 per revision for orchestration, storage, and
delivery. Revision also assumes sixty seconds of non-model parsing and runtime
on a warm Cloudflare Container with one fully active vCPU, 6 GiB of memory, and
12 GB of disk, about \$0.00215 at the published resource
rates~\cite{cfcontainerspricing}. The subtotal excludes payment and chain fees,
refund reserves, training, evaluation, external checks, and low-volume fixed
costs. An idle Container of the standard-2 to standard-4 sizes costs
\$0.057--\$0.113 per hour, so idle capacity and cold starts can dominate per-job cost at low volume. Private-adapter
pricing remains uncertain beyond Cloudflare's advertised free LoRA beta; a
base-model tariff is only a comparison point~\cite{cfloras}.

The dated comparison uses Pangram's standalone realtime rate of \$0.50 per
1,000 words and its advertised bulk rate of \$0.40, against the proposed
\$0.02~\cite{pangrampricing,pangrambulk}. Subscriptions, negotiated prices,
per-item rounding, and differing detection performance limit this comparison.
Prices must be checked again when an offer is published. At
\(M_{\mathrm{price}}=25\), raw-text detection prices at about \$0.007--\$0.011,
at the bottom of the target range; a \$0.02 price corresponds to
\(M_{\mathrm{price}}\) of about 47--70. Price is an independent product
hypothesis, not evidence of comparable accuracy.

Calibration runs the whole admitted pipeline on authorized workloads at
realistic concurrency. Record every attempt's opaque identifier, model,
billed account, usage, latency, outcome, and tariff version without text.
Reconcile usage with provider charges, separating cash charges, promotional
credits, and undiscounted estimates. Divide the cost of all attempts by
successful deliveries, and report failures and cost percentiles separately.
Measure miner and operator expenses separately before selecting \(M_{\mathrm{price}}\), workload
limits, or a paid-use emission share.

\subsection{Checkout and the burn policy}

The planned funding routes are:
\[
\begin{aligned}
\text{USD card payment} &\longrightarrow \text{TAO}
  \longrightarrow \text{alpha} \longrightarrow \text{burn},\\
\text{TAO payment} &\longrightarrow \text{alpha}
  \longrightarrow \text{burn},\\
\text{alpha payment} &\longrightarrow \text{burn}.
\end{aligned}
\]
Coinbase Onramp is the selected candidate for card-to-native-TAO funding.
TAO.com's documented integration supports direct TAO delivery, but Slop
Ninja's access, supported card types and jurisdictions, recipient wallets,
limits, holds, and refunds still need qualification~\cite{coinbaseonramp}.
Credit-card support cannot be inferred from a documented debit-card flow.
This intended route requires neither Ethereum nor a separate bridge. Solana
is the preferred entry network for additional crypto inputs, pending a
qualified route to native TAO.

The target is automated checkout. Onramp delivery, conversion, service delivery,
and burning remain separate events with independent failure modes. Identity
checks, finality, and liquidity can delay them. A card authorization or an
expected exchange fill cannot fund execution; the adapter must confirm arrival
at the authorized funding destination. Provider-hosted forms handle card
details. Payment systems receive opaque job references, not customer documents
or author samples. At cent-scale prices, minimum purchases and fees may require
a funded balance spanning multiple jobs; that product rule remains open.

All inference principal follows the selected burn route. The operator funds
processing, conversion, and chain fees separately, including any shortfall
when a provider deducts fees from its output. Customers bear their own wallet
submission fees unless the offer explicitly sponsors them. Optional external
detector checks have a separate disclosed budget and require authorization to
send the text. Their charges neither reduce inference principal nor earn
additional inference credit.

A quote fixes the original currency and amount, conversion conditions, and
slippage limits. It records an estimated alpha quantity; the final record uses
actual fills and actual alpha destroyed. Native alpha's transfer and staking
semantics, deployment permissions, and burn operations must be qualified
against the runtime. Burning existing alpha does not itself buy TAO, and
recycling is a different operation~\cite{stakeburn,alphaburnimplementation}.
If the deployed burn call is owner-only, a subnet lease must provide an
authorized owner execution path or another verified route before sales open.
Operating the service does not itself grant that permission. This is separate
from the lease's required path for owner-controlled hyperparameter changes
(Section~\ref{participants-and-the-two-kinds-of-work}).
An address transfer or treasury allocation
alone cannot substantiate a burn claim.

Funds remain refundable custody until valid delivery and the published
pre-burn refund window have passed. Conversion may occur earlier if that
custody and its refund liability remain fully funded. Burn batches allocate
each job's principal uniquely and reconcile deposits, conversion fills,
custody, refunds, and finalized burns. Using treasury inventory temporarily
does not replace the promised purchase-and-burn accounting.

\subsection{Reservation, execution, and delivery}\label{reservation-and-delivery}
\label{application-and-agent-integration}

The API and an MCP integration expose quote, reserve, submit, status, and
result operations~\cite{mcp2026}. A quote binds the accepted pipeline and cost
profile, price and currency, input limits, attempt budget, acceptance and
delivery cutoffs, retention, and refund terms. Size and service requirements
suffice for discovery. The document goes to Slop Ninja only after reservation;
later \textit{B} miners receive no customer text through the discovery process.

\begin{enumerate}
\tightlist
\item Reserve confirmed funds, refund liquidity, and capacity durably before
dispatch. A customer idempotency key identifies one job. An expired or
exhausted offer cannot be reserved, and repricing cannot alter an existing
reservation. Cross-ledger operations must reconcile before execution.
\item Validate the submitted document, format, size, and customer output key
before acceptance. Construct the admitted input, canonical text under the
launch profile, and run the selected \textit{A} artifact. Check complete input
coverage, schema validity, probabilities, share/confidence consistency, and
attribution. An input beyond the artifact's length limit, or under a later CFG
profile an unresolved parse or exhausted ambiguity budget, produces an explicit
failure or a separately quoted diagnostic result.
\item For later \textit{B}, execute the bounded proposal and ranking procedure
and check the assembled document for preservation, readability, requested
tone, and any qualified author-fit objective. Local edit acceptance does not
establish whole-document success. Customer-requested Pangram checks require
separate disclosure authorization and funding.
\item Encrypt the result to the customer's output key, retain it in private
storage, and issue an operator-signed receipt. Retrieval must remain available
for the quoted period independently of a serving miner. Complete the required
delivery checks before authorizing burn or inference credit.
\item After the pre-burn window, settle the principal's burn allocation once.
Retain content-free accounting, return excess funds, and expire document data,
derived artifacts, and caches according to the disclosed retention policy,
including backup retention. Training reuse requires separate consent.
\end{enumerate}

The exact delivery predicate is a paid-launch prerequisite. It must bind a
durably stored encrypted object, the authorized customer/key, verified digest
and length, working retrieval, completed validations, and timely delivery.
Whether customer acknowledgment is required, and how acceptance interacts
with the refund window, remains to be fixed in the customer terms. A response
hash alone cannot authorize an irreversible burn. The offer buys a bounded
computation and retrievable result; subjective quality guarantees or additional
revision allowances require explicit terms.

Transactional Durable Object ledgers preserve reservations and transitions.
Queues carry job identifiers; private content stays in the operator account.
Persist each attempt's identity and maximum budget before dispatch. A timed-out
provider call may have executed and incurred charges, so mark its outcome
uncertain and reconcile before a bounded retry. Customer charges, burns, and
credits remain idempotent even when provider execution repeats. Restarts do
not extend accepted customer cutoffs. A cohort pause under
Section~\ref{cohorts} preserves research work but does not suspend delivery,
retrieval, or refund obligations.

Switching to another model or later \textit{B} provider requires either prior
customer authorization that fixes replacement eligibility and budget, or a
fresh quote and acceptance. A replacement cannot inherit the reserved model's
benchmark claim. Retain the original attempt and its outcome when a qualified
fallback completes the service.

\subsection{Receipts and refund liabilities}\label{receipt-certified-payment}
\label{payment-without-disclosing-the-job}

A completion receipt binds the protocol and chain domain, job and acceptance,
salted input commitment, customer output-key identity, enrolled model and
pipeline, attempt-ledger commitment, encrypted result digest and size,
deadline, and signer version. Linked payment records bind the quote,
original-currency amount, conversion history, and refund state. Burn records
bind finalized transactions and each job's allocation. Canonical encodings,
signature verification, key rotation, and replay protection must be fixed
before these records authorize funds.

Each inference attempt also has a content-free R2 receipt included in a certified
cohort package. Section~\ref{cohort-receipt-retention} defines its index, validator
verification, and deletion after two subsequent cohorts complete. That rolling
window applies to detailed inference receipts, separately from customer delivery
terms, compact burn/accounting records, and unresolved refund obligations.

Content-free commitments may be anchored on-chain, with inclusion proofs for
batches. Customer ciphertext, billing identity, report URLs, and raw document
hashes stay off-chain. Use fresh secret salts for document commitments. The
receipt states the operator's assertion about execution; it supplies neither
a Cloudflare signature nor independent validator certification. Result
availability remains an operator obligation. Later \textit{B} receipts report
the qualified identity evidence without claiming stronger historical model
immutability than the provider supports.

Rejection before acceptance releases funds and capacity. Failure to deliver
the agreed service requires the published refund; failed attempts still cost
the billed account holder and earn no inference credit. Valid delivery after
bounded retries keeps the original price. Cancellation after acceptance follows
the fixed terms. A conversion or burn delay retains pending credit and follows
its fixed cutoff. Reconcile finalized chain inclusion before deciding that
a pending transaction failed; a local timeout cannot both return principal
and leave an unresolved burn executable.

The proposed service refund restores the original amount in the original
currency: USD through a qualified return path, TAO to the bound refund address,
or alpha to its qualified position. Conversion losses and consumed provider
fees fall to the operator. Card checkout cannot promise this until an actual
USD return path is qualified; an onramp's completed crypto purchase may be
irreversible. Custodial intermediate alpha is reconciled into the refund,
not returned as an additional payment.

After a burn, a refund is a new operator-funded payment. Hold reserves for
retention obligations, service claims, and uncovered processor disputes;
reconcile chargebacks and manual refunds to avoid duplicate returns. Already
paid native emissions cannot be assumed recoverable. Stop new admissions when
required reserves are unavailable.

\subsection{Who pays, and what paid use earns}\label{revenue-and-costs}

Miners fund model preparation and training. Slop Ninja pays to serve public
\textit{A} artifacts, including their Container execution, storage,
coordination, and delivery. Later \textit{B} miners pay Cloudflare for their own customer
inference and bounded mandatory evaluation calls. They absorb permitted
failures and retries as part of participation. Their required candidate
Pangram reports use prepaid balances and a fixed budget; those deposits are
separate from burned customer principal. Unused, uncommitted balances remain
withdrawable.

The operator funds shared evaluation sets, source baselines,
authoritative \textit{A} execution, and independent review through named budgets.
Validators fund their own record verification and native weight submissions;
shared evaluation and service never depend on their funding or approval.
The operator also pays conversion and burn fees, refunds, and its other service
liabilities. Access to \texttt{rtao-rpc} has no service charge; transaction fees
and the caller's runtime costs remain. Insufficient evaluation funding reduces
new assignments rather than evidence requirements or accepted miner budgets.

These expenses require spendable operating funds. Customer principal is burned,
so a high markup supplies no cash for invoices. Capital or explicitly allocated
and realized owner proceeds may fund operations; future emissions are not
reserves. Native owner emissions accrue under the actual ownership or lease
arrangement and do not automatically fund Slop Ninja or an investor
\cite{mechanisms}. Report operating cash, purchased and burned alpha, miner
emissions, and separately priced enterprise services separately.

Section~\ref{benchmark-scoring-and-mechanism-mapping} defines the proposed
quality and paid-use allocations. Inference credit requires an admitted version,
acceptable delivery, and finalized allocation of paid principal to a burn.
Pending burns retain pending credit. Failed attempts, free traffic, mandatory
evaluations, duplicate receipts, and extra tokens earn none. Work units come
from a frozen workload schedule. A published \textit{A} artifact hosted by the
operator also requires an explicit choice about model-author credit versus
serving-cost credit before its calls enter a paid-use pool.

Assign usage to the epoch containing finalized burn inclusion plus a fixed
protocol lag. Publish that lag, the applicable emission-share schedule,
burn cutoff, and batch policy before admission. The operator cannot select
a favorable reward epoch after observing its pool. Delayed or late burns
follow the quoted resolution rule without creating duplicate principal
dispositions or retroactive work credit.

Section~\ref{token-economics} develops the burn hypothesis, cost-recovery
conditions, and self-purchase incentives. Its sensitivity analysis separates
miner participation from the operator's funding and reserve requirements.
